\chapter{Boolean Analysis — Razborov Smolensky}
%%% generated by the blueprint pipeline from TCSlib.BooleanAnalysis.RazborovSmolensky
%%%%%%%%%%%%%%%%%%%%%%%%%%%%%%%%%%%%%%%%%%%%%%%%%%%%%%%%%%%%%%%%%%%%%%%%%%%%%%%
\section{Overview}

This module completes the Razborov--Smolensky lower bound by transferring the root-cube
inapproximability result back to the Boolean $\mathrm{MOD}_q$ function.  Padding by $q-1$
extra input bits recovers every residue class of the Hamming weight from the zero-residue
$\mathrm{MOD}_q$ predicate; an affine change of variables converts a low-degree approximant
of padded $\mathrm{MOD}_q$ over $\bbz/p$ into low-degree approximants of the residue
indicators on the cube $\{1,\omega\}^n$ over $\bbf_{p^{q-1}}$, and a union bound recombines
these into an approximant of the root product $\prod_i x_i$.  Combined with the counting
lower bound for the root product and the circuit-approximation theorem, this yields
Smolensky's low-degree lower bound for $\mathrm{MOD}_q$ and the quantitative theorem
$\mathrm{MOD}_q \notin \mathrm{AC}^0[p]$.

%%%%%%%%%%%%%%%%%%%%%%%%%%%%%%%%%%%%%%%%%%%%%%%%%%%%%%%%%%%%%%%%%%%%%%%%%%%%%%%
\section{Declarations}

\begin{definition}[Boolean embedding into the root cube]
\lean{ACP.boolToRootCube}
\label{ACP.boolToRootCube}
\leanok
\uses{ACP.rootCube}
Let $K$ be a field and $\omega \in K$.  A Boolean input $x \in \{0,1\}^n$ is embedded as
the point of the cube $\{1,\omega\}^n$ whose $i$-th coordinate is $1$ if $x_i = 0$ and
$\omega$ if $x_i = 1$.
\end{definition}

\begin{theorem}[Root product of an embedded Boolean input]
\lean{ACP.boolToRootCube_rootProduct_eq_weightPow}
\label{ACP.boolToRootCube_rootProduct_eq_weightPow}
\leanok
\uses{ACP.boolToRootCube}
\difficulty{4}
Let $K$ be a field, $\omega \in K$, and let $x \in \{0,1\}^n$ be a Boolean input.  The
product of the coordinates of the point of $\{1,\omega\}^n$ obtained from $x$ by the
embedding $0 \mapsto 1$, $1 \mapsto \omega$ equals $\omega^{\abs{x}}$, where
$\abs{x} = \sum_i x_i$ is the Hamming weight of $x$.
\end{theorem}

\begin{definition}[$\omega$-weight of a root-cube point]
\lean{ACP.rootResidueWeight}
\label{ACP.rootResidueWeight}
\leanok
\uses{ACP.rootCube}
Let $K$ be a field and $\omega \in K$.  The $\omega$-weight of a point
$x \in \{1,\omega\}^n$ is the number of indices $i$ with $x_i = \omega$.  Under the
Boolean embedding $0 \mapsto 1$, $1 \mapsto \omega$, this is the Boolean Hamming weight.
\end{definition}

\begin{definition}[Residue indicator on the root cube]
\lean{ACP.rootResidueIndicator}
\label{ACP.rootResidueIndicator}
\leanok
\uses{ACP.rootCube, ACP.rootResidueWeight}
Let $K$ be a field, $\omega \in K$, and let $q$ be a positive natural number.  For a
residue $r \in \{0,\dots,q-1\}$, the residue-$r$ indicator is the $K$-valued function on
the cube $\{1,\omega\}^n$ that takes the value $1$ at $x$ when the $\omega$-weight of $x$
is congruent to $r$ modulo $q$, and $0$ otherwise.
\end{definition}

\begin{definition}[Boolean Hamming weight]
\lean{ACP.boolWeight}
\label{ACP.boolWeight}
\leanok
The Hamming weight of a Boolean input $x \in \{0,1\}^n$, i.e.\ the natural number
$\sum_{i} x_i$.
\end{definition}

\begin{definition}[Boolean residue indicator]
\lean{ACP.boolResidueIndicator}
\label{ACP.boolResidueIndicator}
\leanok
\uses{ACP.boolWeight}
Let $K$ be a field and let $q$ be a positive natural number.  For a residue
$r \in \{0,\dots,q-1\}$, the Boolean residue-$r$ indicator is the $K$-valued function on
Boolean inputs $x \in \{0,1\}^n$ that is $1$ when the Hamming weight of $x$ is congruent
to $r$ modulo $q$, and $0$ otherwise.
\end{definition}

\begin{definition}[Decoded bit of a root-cube coordinate]
\lean{ACP.rootCubeBit}
\label{ACP.rootCubeBit}
\leanok
\uses{ACP.rootCube}
Let $K$ be a field and $\omega \in K$.  The $i$-th decoded bit of a point
$x \in \{1,\omega\}^n$ is the Boolean value that is $0$ if $x_i = 1$ and $1$ otherwise.
When $\omega \neq 1$ this inverts the Boolean embedding, sending the coordinate value $1$
to the bit $0$ and the value $\omega$ to the bit $1$.
\end{definition}

\begin{definition}[Number of padding ones for a residue]
\lean{ACP.residuePadOnes}
\label{ACP.residuePadOnes}
\leanok
Let $q$ be a positive natural number.  For a residue $r \in \{0,\dots,q-1\}$, the number
of padding $1$-bits needed to shift Hamming-weight residue $r$ to residue $0$ modulo $q$,
namely $(q - r) \bmod q$: this is $0$ when $r = 0$ and $q - r$ when $r > 0$.
\end{definition}

\begin{definition}[Padded Boolean input of a root-cube point]
\lean{ACP.paddedResidueInput}
\label{ACP.paddedResidueInput}
\leanok
\uses{ACP.residuePadOnes, ACP.rootCube, ACP.rootCubeBit}
Let $K$ be a field, $\omega \in K$, let $q$ be a positive natural number, let
$r \in \{0,\dots,q-1\}$, and let $x \in \{1,\omega\}^n$.  The associated padded Boolean
input on $n + (q-1)$ bits has as its first $n$ bits the decoded bits of $x$ (bit $0$
where $x_i = 1$, bit $1$ elsewhere), while its final $q-1$ bits consist of
$(q - r) \bmod q$ ones followed by zeros.
\end{definition}

\begin{definition}[Affine Boolean decoder polynomial]
\lean{ACP.rootAffineBoolPoly}
\label{ACP.rootAffineBoolPoly}
\leanok
Let $K$ be a field and $\omega \in K$.  For a variable index $i \in \{1,\dots,n\}$, the
affine polynomial $(\omega - 1)^{-1}\,(X_i - 1)$ in $n$ variables over $K$.  When
$\omega \neq 1$ it decodes a root-cube coordinate as a Boolean value, sending
$1 \mapsto 0$ and $\omega \mapsto 1$.
\end{definition}

\begin{theorem}[Degree bound for the affine decoder]
\lean{ACP.rootAffineBoolPoly_totalDegree_le_one}
\label{ACP.rootAffineBoolPoly_totalDegree_le_one}
\leanok
\uses{ACP.rootAffineBoolPoly}
\difficulty{3}
Let $K$ be a field, $\omega \in K$, and let $i$ be a variable index.  The affine decoder
polynomial $(\omega - 1)^{-1}\,(X_i - 1)$ has total degree at most $1$.
\end{theorem}

\begin{definition}[Residue substitution for padded variables]
\lean{ACP.paddedResidueSubst}
\label{ACP.paddedResidueSubst}
\leanok
\uses{ACP.ModqField, ACP.residuePadOnes, ACP.rootAffineBoolPoly}
Let $p$ and $q$ be primes, let $K = \bbf_{p^{q-1}}$, let $\omega \in K$, and let
$r \in \{0,\dots,q-1\}$.  The residue-$r$ substitution assigns to each of the $n + (q-1)$
variables a polynomial in $n$ variables over $K$: a variable $X_j$ with $j < n$ is sent
to the affine decoder $(\omega - 1)^{-1}\,(X_j - 1)$, while the $k$-th padding variable
(with $k = j - n$) is sent to the constant $1$ if $k < (q - r) \bmod q$ and to the
constant $0$ otherwise.
\end{definition}

\begin{definition}[Residue restriction of a padded polynomial]
\lean{ACP.paddedResiduePolynomial}
\label{ACP.paddedResiduePolynomial}
\leanok
\uses{ACP.ModqField, ACP.paddedResidueSubst}
Let $p$ and $q$ be primes, let $K = \bbf_{p^{q-1}}$, let $\omega \in K$, let
$r \in \{0,\dots,q-1\}$, and let $P$ be a polynomial in $n + (q-1)$ variables over
$\bbz/p$.  The residue-$r$ restriction of $P$ is the polynomial in $n$ variables over $K$
obtained by extending the coefficients of $P$ along the canonical embedding
$\bbz/p \to K$ and then substituting each variable according to the residue-$r$
substitution: the first $n$ variables become the affine root-cube decoders and the
padding variables are fixed to the constants encoding residue $r$.
\end{definition}

\begin{theorem}[Residue restriction does not increase degree]
\lean{ACP.paddedResiduePolynomial_totalDegree_le}
\label{ACP.paddedResiduePolynomial_totalDegree_le}
\leanok
\uses{ACP.ModqField, ACP.paddedResiduePolynomial, ACP.paddedResidueSubst, ACP.residuePadOnes, ACP.rootAffineBoolPoly_totalDegree_le_one}
\difficulty{5}
Let $p$ and $q$ be primes, let $\omega \in \bbf_{p^{q-1}}$, and let
$r \in \{0,\dots,q-1\}$.  For every polynomial $P$ in $n + (q-1)$ variables over
$\bbz/p$, the residue-$r$ restriction of $P$ (a polynomial in $n$ variables over
$\bbf_{p^{q-1}}$) has total degree at most the total degree of $P$.
\end{theorem}

\begin{theorem}[Evaluation of the residue restriction]
\lean{ACP.paddedResiduePolynomial_eval_eq_map_eval}
\label{ACP.paddedResiduePolynomial_eval_eq_map_eval}
\leanok
\uses{ACP.ModqField, ACP.boolInput, ACP.paddedResidueInput, ACP.paddedResiduePolynomial, ACP.paddedResidueSubst, ACP.residuePadOnes, ACP.rootAffineBoolPoly, ACP.rootCube, ACP.rootCubeBit}
\difficulty{5}
Let $p$ and $q$ be primes, let $\omega \in \bbf_{p^{q-1}}$ with $\omega \neq 1$, let
$r \in \{0,\dots,q-1\}$, let $P$ be a polynomial in $n + (q-1)$ variables over $\bbz/p$,
and let $x \in \{1,\omega\}^n$.  Evaluating the residue-$r$ restriction of $P$ at $x$
gives the same value as evaluating $P$ at the padded Boolean input associated to
$(x, r)$ (with bits cast into $\bbz/p$) and then mapping the result along the embedding
$\bbz/p \to \bbf_{p^{q-1}}$.
\end{theorem}

\begin{theorem}[$\mathrm{MOD}_q$ of the padded input is the residue indicator]
\lean{ACP.paddedResidueInput_modQTarget_eq_residueIndicator}
\label{ACP.paddedResidueInput_modQTarget_eq_residueIndicator}
\leanok
\uses{ACP.ModqField, ACP.modQTarget, ACP.paddedResidueInput, ACP.residuePadOnes, ACP.rootCube, ACP.rootCubeBit, ACP.rootResidueIndicator, ACP.rootResidueWeight}
\difficulty{6}
Let $p$ and $q$ be primes, let $\omega \in \bbf_{p^{q-1}}$ with $\omega \neq 1$, let
$r \in \{0,\dots,q-1\}$, and let $x \in \{1,\omega\}^n$.  Write $\mathrm{MOD}_q$ for the
$\bbz/p$-valued function on Boolean inputs that is $1$ when the Hamming weight of its
input is divisible by $q$ and $0$ otherwise.  Then the image in $\bbf_{p^{q-1}}$ of
$\mathrm{MOD}_q$ evaluated at the padded Boolean input associated to $(x, r)$ equals the
residue-$r$ indicator of $x$: it is $1$ exactly when the $\omega$-weight of $x$ is
congruent to $r$ modulo $q$, and $0$ otherwise.
\end{theorem}

\begin{theorem}[Injectivity of the padded-input map]
\lean{ACP.paddedResidueInput_injective}
\label{ACP.paddedResidueInput_injective}
\leanok
\uses{ACP.ModqField, ACP.paddedResidueInput, ACP.rootCubeBit}
\difficulty{4}
Let $p$ and $q$ be primes, let $\omega \in \bbf_{p^{q-1}}$ with $\omega \neq 1$, and fix
a residue $r \in \{0,\dots,q-1\}$.  The map sending a root-cube point
$x \in \{1,\omega\}^n$ to its padded Boolean input on $n + (q-1)$ bits is injective.
\end{theorem}

\begin{theorem}[Residue approximants from one padded $\mathrm{MOD}_q$ approximant]
\lean{ACP.padded_modQ_approx_gives_root_residue_approximants}
\label{ACP.padded_modQ_approx_gives_root_residue_approximants}
\leanok
\uses{ACP.ModqField, ACP.badInputCount, ACP.boolInput, ACP.modQTarget, ACP.paddedResidueInput, ACP.paddedResidueInput_injective, ACP.paddedResidueInput_modQTarget_eq_residueIndicator, ACP.paddedResiduePolynomial, ACP.paddedResiduePolynomial_eval_eq_map_eval, ACP.paddedResiduePolynomial_totalDegree_le, ACP.rootCube, ACP.rootCubeBadCount, ACP.rootResidueIndicator}
\difficulty{5}
Let $p$ and $q$ be distinct primes and let $\omega \in \bbf_{p^{q-1}}$ satisfy
$\omega^q = 1$ and $\omega \neq 1$.  Suppose $P$ is a polynomial in $n + (q-1)$
variables over $\bbz/p$ of total degree at most $d$ that disagrees with the
$\mathrm{MOD}_q$ function (the indicator, in $\bbz/p$, that the Hamming weight is
divisible by $q$) on fewer than $E$ Boolean inputs of length $n + (q-1)$.  Then there is
a family $(R_r)_{r \in \{0,\dots,q-1\}}$ of polynomials in $n$ variables over
$\bbf_{p^{q-1}}$ such that each $R_r$ has total degree at most $d$ and disagrees with the
residue-$r$ indicator on at most $E$ points of the cube $\{1,\omega\}^n$.
\end{theorem}

\begin{theorem}[Residue expansion of the root product]
\lean{ACP.rootProduct_eq_residueIndicator_sum}
\label{ACP.rootProduct_eq_residueIndicator_sum}
\leanok
\uses{ACP.rootCube, ACP.rootResidueIndicator, ACP.rootResidueWeight}
\difficulty{5}
Let $K$ be a field, let $q$ be a prime, and let $\omega \in K$ satisfy $\omega^q = 1$.
For every point $x \in \{1,\omega\}^n$,
\[
  \prod_{i=1}^{n} x_i
    = \sum_{r=0}^{q-1} \omega^{r} \cdot \ind{w(x) \equiv r \pmod q},
\]
where $w(x)$ denotes the number of coordinates of $x$ equal to $\omega$.
\end{theorem}

\begin{theorem}[Union bound for residue recombination]
\lean{ACP.rootCubeBadCount_residueCombination_le_sum}
\label{ACP.rootCubeBadCount_residueCombination_le_sum}
\leanok
\uses{ACP.rootCube, ACP.rootCubeBadCount, ACP.rootProduct_eq_residueIndicator_sum, ACP.rootResidueIndicator}
\difficulty{5}
Let $K$ be a finite field, let $q$ be a prime, and let $\omega \in K$ satisfy
$\omega^q = 1$.  For any family $(R_r)_{r \in \{0,\dots,q-1\}}$ of polynomials in $n$
variables over $K$, the number of points of $\{1,\omega\}^n$ at which the combined
polynomial $\sum_{r} \omega^{r} R_r$ differs from the root product $\prod_i x_i$ is at
most the sum over $r$ of the number of points at which $R_r$ differs from the
residue-$r$ indicator.
\end{theorem}

\begin{theorem}[Combining residue approximants into a root-product approximant]
\lean{ACP.root_residue_approximants_combine_to_rootProduct}
\label{ACP.root_residue_approximants_combine_to_rootProduct}
\leanok
\uses{ACP.rootCube, ACP.rootCubeBadCount, ACP.rootCubeBadCount_residueCombination_le_sum, ACP.rootResidueIndicator}
\difficulty{4}
Let $K$ be a finite field, let $q$ be a prime, and let $\omega \in K$ satisfy
$\omega^q = 1$ and $\omega \neq 1$.  Given polynomials $R_r$ in $n$ variables over $K$,
one for each residue $r \in \{0,\dots,q-1\}$, each of total degree at most $d$ and each
disagreeing with the residue-$r$ indicator on at most $E$ points of $\{1,\omega\}^n$,
there exists a polynomial $Q$ in $n$ variables over $K$ of total degree at most $d$ that
disagrees with the root product $x \mapsto \prod_i x_i$ on at most $q \cdot E$ points of
the cube.
\end{theorem}

\begin{theorem}[Sharp Boolean-to-root-product transfer]
\lean{ACP.padded_modQ_approx_transfers_to_rootProduct_qE}
\label{ACP.padded_modQ_approx_transfers_to_rootProduct_qE}
\leanok
\uses{ACP.ModqField, ACP.badInputCount, ACP.modQTarget, ACP.padded_modQ_approx_gives_root_residue_approximants, ACP.rootCube, ACP.rootCubeBadCount, ACP.root_residue_approximants_combine_to_rootProduct}
\difficulty{2}
Let $p$ and $q$ be distinct primes and let $\omega \in \bbf_{p^{q-1}}$ satisfy
$\omega^q = 1$ and $\omega \neq 1$.  If $P$ is a polynomial in $n + (q-1)$ variables
over $\bbz/p$ of total degree at most $d$ that disagrees with the $\mathrm{MOD}_q$
function (the indicator that the Hamming weight is divisible by $q$) on fewer than $E$
Boolean inputs, then there exists a polynomial $Q$ in $n$ variables over
$\bbf_{p^{q-1}}$ of total degree at most $d$ that disagrees with the root product
$x \mapsto \prod_i x_i$ on at most $q \cdot E$ points of the cube $\{1,\omega\}^n$.
\end{theorem}

\begin{theorem}[Boolean-to-root-product transfer]
\lean{ACP.padded_modQ_approx_transfers_to_rootProduct}
\label{ACP.padded_modQ_approx_transfers_to_rootProduct}
\leanok
\uses{ACP.ModqField, ACP.badInputCount, ACP.modQTarget, ACP.padded_modQ_approx_transfers_to_rootProduct_qE, ACP.rootCube, ACP.rootCubeBadCount}
\difficulty{2}
Let $p$ and $q$ be distinct primes, let $\omega \in \bbf_{p^{q-1}}$ satisfy
$\omega^q = 1$ and $\omega \neq 1$, and let $E$ and $e$ be natural numbers with
$q \cdot E \le e$.  If $P$ is a polynomial in $n + (q-1)$ variables over $\bbz/p$ of
total degree at most $d$ that disagrees with the $\mathrm{MOD}_q$ function (the
indicator that the Hamming weight is divisible by $q$) on fewer than $E$ Boolean inputs,
then there exists a polynomial $Q$ in $n$ variables over $\bbf_{p^{q-1}}$ of total
degree at most $d$ that disagrees with the root product $x \mapsto \prod_i x_i$ on at
most $e$ points of the cube $\{1,\omega\}^n$.
\end{theorem}

\begin{theorem}[Boolean lower bound from root-product inapproximability]
\lean{ACP.modQ_lowDegreeBadCountLB_from_rootProduct}
\label{ACP.modQ_lowDegreeBadCountLB_from_rootProduct}
\leanok
\uses{ACP.LowDegreeBadCountLB, ACP.ModqField, ACP.badInputCount, ACP.modQTarget, ACP.padded_modQ_approx_transfers_to_rootProduct, ACP.rootCube, ACP.rootCubeBadCount}
\difficulty{3}
Let $p$ and $q$ be distinct primes, let $\omega \in \bbf_{p^{q-1}}$ satisfy
$\omega^q = 1$ and $\omega \neq 1$, and let $d, E, e$ be natural numbers with
$q \cdot E \le e$.  Suppose no polynomial in $n$ variables over $\bbf_{p^{q-1}}$ of
total degree at most $d$ disagrees with the root product $x \mapsto \prod_i x_i$ on at
most $e$ points of $\{1,\omega\}^n$.  Then every polynomial in $n + (q-1)$ variables
over $\bbz/p$ of total degree at most $d$ disagrees with the $\mathrm{MOD}_q$ function
(the indicator that the Hamming weight is divisible by $q$) on at least $E$ Boolean
inputs.
\end{theorem}

\begin{theorem}[Smolensky's low-degree lower bound for $\mathrm{MOD}_q$]
\lean{ACP.smolensky_modQ_lowDegreeBadCountLB}
\label{ACP.smolensky_modQ_lowDegreeBadCountLB}
\leanok
\uses{ACP.LowDegreeBadCountLB, ACP.LowDegreeSupport, ACP.ModqField, ACP.exists_nontrivial_qth_root_modqField, ACP.modQTarget, ACP.modQ_lowDegreeBadCountLB_from_rootProduct, ACP.no_low_degree_rootProd_approx_concrete, ACP.rootCube, ACP.rootCubeBadCount}
\difficulty{3}
Let $p$ and $q$ be distinct primes, write $K = \bbf_{p^{q-1}}$, and let $d, E, e, B$ be
natural numbers with $q \cdot E \le e$.  Assume the counting bounds
\[
  \sum_{t=0}^{e} \binom{2^n}{t}\,\abs{K}^{t} \le B
  \qquad\text{and}\qquad
  \abs{K}^{2^n} > N \cdot B,
\]
where $N$ is the number of $K$-valued functions on the collection of subsets
$S \subseteq \{1,\dots,n\}$ with $\abs{S} \le \lfloor n/2 \rfloor + d$.  Then every
polynomial in $n + (q-1)$ variables over $\bbz/p$ of total degree at most $d$ disagrees
with the $\mathrm{MOD}_q$ function (the indicator that the Hamming weight is divisible
by $q$) on at least $E$ Boolean inputs.
\end{theorem}

\begin{theorem}[Quantitative $\mathrm{MOD}_q \notin \mathrm{AC}^0[p]$]
\lean{ACP.MODq_notin_AC0p_quantitative}
\label{ACP.MODq_notin_AC0p_quantitative}
\leanok
\uses{ACP.ACp_GateOps, ACP.FeedForward, ACP.FeedForward.eval₁, ACP.FeedForward.onlyUsesGates, ACP.FeedForward.size, ACP.LowDegreeBadCountLB, ACP.LowDegreeSupport, ACP.ModqField, ACP.circuitDegreeBound, ACP.modGateOp, ACP.modQTarget, ACP.size_lower_bound_from_relative_badCountLB, ACP.smolensky_modQ_lowDegreeBadCountLB}
\difficulty{2}
Let $p$ and $q$ be distinct primes and write $K = \bbf_{p^{q-1}}$.  Let $F$ be a
feed-forward Boolean circuit on $n + (q-1)$ input bits with a single output and finitely
many nodes in each layer, using only $\mathrm{AC}^0[p]$ gate operations, and suppose
that $F$ computes the $\mathrm{MOD}_q$ function: on every input, $F$ outputs $1$ exactly
when the Hamming weight of the input is divisible by $q$.  Let $\delta, \ell, e, B$ be
natural numbers such that $q \cdot \bigl(\delta \cdot 2^{\,n+(q-1)}\bigr) \le e$, such
that $\sum_{t=0}^{e} \binom{2^n}{t}\,\abs{K}^{t} \le B$, and such that
$\abs{K}^{2^n} > N \cdot B$, where $N$ is the number of $K$-valued functions on the
collection of subsets $S \subseteq \{1,\dots,n\}$ with
$\abs{S} \le \lfloor n/2 \rfloor + \bigl((p-1)\ell\bigr)^{\mathrm{depth}(F)}$.  Then the
size of $F$ (its number of non-input gates) satisfies
$\delta \cdot 2^{\ell} \le \mathrm{size}(F)$.
\end{theorem}
